If $\mathfrak p$ is not in $V(\mathfrak a)$, then $\mathfrak p$ does not contain some element $f\in\mathfrak a$ and $D(f)$ is a neighbourhood around $\mathfrak p$ in the complement of $V(\mathfrak a)$, so $D(f)$ form a basis of open sets.

\textbf{(i)} Prime $\mathfrak p$ does not contain $f$ and $g$ iff it does not contain $fg$.

\textbf{(ii)} If $f$ is in all $\mathfrak p$ then $f$ in nilradical.

\textbf{(iii)} If $f$ is in no prime ideal then it is in no maximal ideal, so it is a unit.

\textbf{(iv)} $\sqrt{(f)}=\sqrt{(g)}$ because they are intersections of elements in $V(f)=V(g)$ by Prop 1.14.

\textbf{(v)} Say $X$ covered by $D(f_i)$ (can decompose arbitrary cover into this form because $D(f)$ form a basis), then $(f_i)_{i\in I}=(1)$ because otherwise the maximal ideal containing $(f_i)_{i\in I}$ is in $X$ but not in any $D(f_i)$. Write $r_1f_1+\cdots+r_nf_n=1$ for some $r_i\in A$, then $\{D(f_i)\}_{i=1}^n$ covers $X$ because a prime can not contain all of $f_1,\ldots,f_n$.

\textbf{(vi)} This is because $X_f=\operatorname{Spec}A_f$. Alternatively, given cover $D(f_i)$, since all primes containing $\mathfrak a$ must contain $f$, we have $f$ is nilpotent in $X/\mathfrak a$ so $f\in\sqrt{\mathfrak a}$. Take $f_1,\ldots,f_n$ such that $(f_1,\ldots,f_n)$ contains a power of $f$, then $D(f_1),\ldots,D(f_n)$ covers $X_f$.

\textbf{(vii)} An open set is a union of $D(f)$ since they form a basis, if it were quasi-compact, then it is a finite union.
